El dual de $(P)$ es:

\[
\begin{aligned}
\min \quad &
\sum_{j=1}^{n} c_j x_j\\[2mm]
\text{s.a.}\quad &
\sum_{j=1}^{n} a_{ij}x_j \ge b_i,
&& i=1,\ldots,m,\\[2mm]
&
\sum_{j=1}^{n} s_{kj}x_j \le q_k,
&& k=1,\ldots,p,\\[2mm]
&
x_j\ge 0,
&& j=1,\ldots,n.
\end{aligned}
\]

La variable $x_j$ representa la cantidad o número de porciones del alimento $j$ que se incluyen en la dieta.

% \begin{itemize}
%     \item La función objetivo minimiza el coste total de la dieta.
%     \item Las restricciones
%     \[
%     \sum_{j=1}^{n} a_{ij}x_j \ge b_i
%     \]
%     garantizan que se cubran los requerimientos mínimos de cada nutriente esencial.
%     \item Las restricciones
%     \[
%     \sum_{j=1}^{n} s_{kj}x_j \le q_k
%     \]
%     garantizan que no se superen los límites máximos permitidos de los componentes críticos.
% \end{itemize}
